\begin{table*}[t]
\centering
\small
\setlength{\tabcolsep}{4pt}
\caption{Claim-to-evidence map for the final KBS framing. The table states which experiments carry each substantive claim and, equally importantly, what each evidence block does not establish. This is the intended reading order for the main and supplementary audits.}
\label{tab:claim-evidence-map}
\maxtablewidth{
\begin{tabular}{@{}p{3.3cm}p{4.6cm}p{4.9cm}p{3.8cm}@{}}
\toprule
Claim & Main evidence & Supported reading & Not established \\
\midrule
The released \texttt{tracer\_adaptive} object is an auditable single-route deployment policy. & \Cref{tab:tracer-route-audit,tab:tracer-route-rules,tab:tracer-rule-rationale,tab:route-signature-provenance,tab:tracer-threshold-intervals,tab:tracer-threshold-lobo} & Route choice follows frozen train/dev predicates and can be inspected as a maintained knowledge artifact rather than as hidden test-time model selection. & A fully formal knowledge model or uniquely identified thresholds on every benchmark family. \\
Retrieval-active TRACER routes matter on the primary event-disjoint setting. & \Cref{tab:atlasv2-event-significance,tab:atlasv2-lopo-family-audit,tab:atlasv2-public-event-extended,tab:clean-mechanism-ablation,tab:atlasv2-lopo-core-ablation} & The released retrieval-active route beats the individual main-table baselines on the small ATLASv2 held-out-family audit and remains competitive over processed-window LOPO folds. & Universal dominance over all fixed families or proof that every calibration/gate component is individually necessary. \\
Case retrieval provides inspectable evidence when the policy keeps retrieval active. & \Cref{tab:decision-trace-examples,fig:warning-score-evidence-alignment,tab:knowledge-grounded-evidence,tab:case-evidence-utility,tab:atlasv2-review-augmentation,tab:alert-stability-audit} & High-score warnings often surface semantically anchored train-only analogs and stable alert episodes, especially on the ATLASv2 held-out-family route. & A human-validated triage gain or an independent top-budget ranking boost from analog labels alone. \\
Simple baselines can be strong on tiny short-prefix splits, but they do not remove the routed-policy case. & \Cref{tab:tabular-baseline-audit,tab:policy-v2-tabular-gate,tab:atlasv2-public-main,tab:aitads-public-main,tab:atlasraw-public-main} & Flat-prefix logistic can be stronger on the tiny ATLASv2 chronology split, whereas a train/dev-only tabular admission check would not select that ex-post row and broader public behavior is less favorable to a universal classical-baseline story. & A claim that the adaptive policy is the best short-prefix scorer on every chronology benchmark. \\
Transparent rank consensus is the stronger score layer when multi-model serving is acceptable. & \Cref{tab:cross-benchmark-rank-ensemble,tab:ensemble-selector-audit,tab:selector-regret-audit} & Calibration-free rank fusion is stronger than the single route on the main prediction caches, and adding the adaptive score gives small extra gains. & A claim that the single-route policy is the final best score aggregator once multi-model serving is allowed. \\
External behavior is competitive but mixed. & \Cref{tab:aitads-public-main,tab:atlasraw-public-main,tab:splunk-attack-data-probe,tab:cross-dataset-transfer} & The frozen family remains competitive across additional public settings and can recover with small target support under some shifts. & Zero-shot universal transfer or external dominance across every public corpus. \\
\bottomrule
\end{tabular}
}
\end{table*}
